\chapter{Development lineage: the two-stage pipeline}\label{app:lineage}
The proposed model was reached through an earlier version, the two-stage
pipeline, which combines a plug-in mean with a per-channel graph GP on the
residual. The final model contains that pipeline as the corner of its
constraint space in which every constraint is switched on
(\S\ref{sec:special-cases}), so the earlier version survives as a nested
special case of the proposed model and as its ablation chain. This appendix
records the per-channel, precision-side formulation that the pipeline used,
since that is the form in which the graph prior is most easily read. The
posteriors of the constrained models, the ablation sequence of the
development, the selection history and the two-stage results of the
downstream tasks form the development record and are kept in the repository
beside the code.\footnote{The development record is the design document
\texttt{docs/graphgp\_first\_design.md} in the repository, together with the
results documents it names. It also records the ablation sequence and the
selection history that led from the per-channel form of this appendix to the
proposed model.} The main line carries the comparison of the two models on
the held-out test set (Table~\ref{tab:headline-conf}).

\section{The per-channel form, precision-side}
With $B$ diagonal and the feature kernels switched off, the joint GP
factorises into three scalar graph GPs, and each of these is more naturally
written through its \emph{precision}, the inverse covariance
$\Qg = K_G^{-1}$ scaled per channel. The two precisions used in the
development sequence are
\begin{align}
  \text{additive:} \quad & \Qg = \lambda I + \eta\, \Lg,
      && \lambda > 0,\;\; \eta \ge 0;
  \label{eq:additive-prec}\\[2pt]
  \text{graph Mat\'ern / SPDE:} \quad & \Qg = \sigma_g^{-2}\,(\kappa^2 I + \Lg)^{\alpha},
      && \sigma_g > 0,\;\; \kappa > 0,\;\; \alpha \in \mathbb{Z}_{\ge 1},
  \label{eq:matern-prec}
\end{align}
where the additive form matches \eqref{eq:additive} with shape
$s = \eta / \lambda$ and scale $1/\lambda$, the scale being carried by
$B_{cc}$ in the normalised parameterisation, and the Mat\'ern form matches
\eqref{eq:matern} in the same way. The additive precision has a direct
reading. The term $\eta\,\Lg$ penalises differences between connected notes,
while the ridge $\lambda I$ pulls the field toward its mean and sets the
overall scale. The choice between precision and kernel is computational
rather than a modelling choice. The precision $\Qg$ is sparse, because each
note couples only to its graph neighbours, which is precisely the Markov
property of a \textbf{Gaussian Markov random field} (GMRF), and this
sparsity makes the per-channel posterior below cheap. The full model, by
contrast, is naturally written on the covariance side, since the
coregionalization and the feature kernels add dense terms, and it is
evaluated spectrally through \eqref{eq:spectral} with one eigendecomposition
per piece.

With a plug-in mean $\mu$ and a diagonal observation precision $P$, equal to
$1/\sigma_e^2$ on observed notes and $0$ on masked notes, the per-channel
posterior is
\begin{equation}
  \Sy \;=\; \bigl(\Qg + P\bigr)^{-1}, \qquad
  m \;=\; \Sy \bigl(\Qg\, \mu + P\, \tilde y\bigr)
  \label{eq:posterior-chan}
\end{equation}
which is the precision-side counterpart of \eqref{eq:posterior}, and to
which \eqref{eq:posterior} reduces under the constraints. The plug-in mean
itself is one of three choices. The first is $\mu = 0$. The second is a
ridge regression on the hand-built features $X_{\mathrm{feat}}$. The third
is the \textbf{music-model mean} $\muLM$, a ridge head $W_{\mathrm{head}}$,
fitted once on a piece-disjoint split, applied to the embeddings of
\S\ref{sec:features},
\begin{equation}
  \mu_{\mathrm{LM},i} = W_{\mathrm{head}}^{\top}\,[\, h_i;\, 1\,],
  \qquad
  y = \muLM + \xi, \qquad \xi \sim \N\!\bigl(0,\, \Qg^{-1}\bigr),
  \label{eq:lm-mean}
\end{equation}
with the graph modelling the residual $y - \muLM$. The hyperparameters
$\theta = (\lambda, \eta, \sigma_e^2)$ are fitted per channel by the same
marginal likelihood \eqref{eq:marglik} restricted to that channel, with the
plug-in mean subtracted from $\tilde y$, with the same noise floor, and with
a per-channel safeguard. The safeguard is a fallback sequence: the marginal
likelihood first, then the calibration-split objective, then the
conservative decoupled fit with $\eta = 0$ and
$\lambda = 1/\widehat{\operatorname{Var}}[\tilde y_o - \mu_o]$. The
per-channel safeguard differs from that of the joint model in one respect.
Its screen tests \emph{recovery} only, that is, the error against the
mean-only fallback, whereas the overconfidence screen of \S\ref{sec:eb}
exists only in the safeguard of the joint model. The two screens are
related, not identical, and the ablation sequence that led from this form to
the proposed model is part of the development record named above.

% =====================================================================
